\documentclass[a4paper,11pt]{article}


\usepackage[T1]{fontenc}
\usepackage[utf8]{inputenc}
\usepackage{graphicx}
\usepackage{xcolor}
\usepackage{mathptmx}


\usepackage{amsmath,amssymb,amsthm,textcomp}
\usepackage{multicol}
\usepackage{tikz}
\usepackage[normalem]{ulem}
\usepackage{bbm} % using for the characteristic function

\usepackage{geometry}
\geometry{left=25mm,right=25mm,%
bindingoffset=0mm, top=20mm,bottom=20mm}
\usepackage{pdfpages} %pdf




%\newtheorem{theorem}{Theorem}
\newtheorem{theorem}{Theorem}[section]
\newtheorem{review}{Review}[section]
\newtheorem{corollary}[theorem]{Corollary}
\newtheorem{lemma}[theorem]{Lemma}
\newtheorem{algorithm}[theorem]{Algorithm}

\theoremstyle{definition}
\newtheorem{definition}[theorem]{Definition}
\newtheorem{example}[theorem]{Example}

\theoremstyle{remark}
\newtheorem*{remark}{Remark}


\linespread{1.3}

\newcommand{\linia}{\rule{\linewidth}{0.5pt}}


% my own titles
\makeatletter
\renewcommand{\maketitle}{
\begin{center}
\vspace{2ex}
{\huge \textsc{\@title}}
\vspace{1ex}
\\
\linia\\
\@author \hfill \@date
\vspace{4ex}
\end{center}
}
\makeatother
%%%

% custom footers and headers
\usepackage{fancyhdr}
\pagestyle{fancy}
\lhead{}
\chead{}
\rhead{}
\lfoot{}
\cfoot{}
\rfoot{\thepage}
\renewcommand{\headrulewidth}{0pt}
\renewcommand{\footrulewidth}{0pt}
%

% code listing settings
\usepackage{listings}
\lstset{
    language=Python,
    basicstyle=\ttfamily\small,
    aboveskip={1.0\baselineskip},
    belowskip={1.0\baselineskip},
    columns=fixed,
    extendedchars=true,
    breaklines=true,
    tabsize=4,
    prebreak=\raisebox{0ex}[0ex][0ex]{\ensuremath{\hookleftarrow}},
    frame=lines,
    showtabs=false,
    showspaces=false,
    showstringspaces=false,
    keywordstyle=\color[rgb]{0.627,0.126,0.941},
    commentstyle=\color[rgb]{0.133,0.545,0.133},
    stringstyle=\color[rgb]{01,0,0},
    numbers=left,
    numberstyle=\small,
    stepnumber=1,
    numbersep=10pt,
    captionpos=t,
    escapeinside={\%*}{*)}
}

%%%----------%%%----------%%%----------%%%----------%%%
\begin{document}
\includepdf[pages = -]{sheet1.pdf}

\section*{Set 1}
When you use experiment to varify some thing, you need to quantify the significance\\
data driven, not model dependant
\subsection*{Aufgabe 2}[Glivenko-Cantalli]
\begin{theorem}
	Given a series $\sum f_n$, if $\| f_n \|_\infty < a_k$ with $\sum a_k$ converges. We know $\sum f_n$ converges uniformly. 
\end{theorem}
\begin{remark}
Not useful since we can't find the Majorant. Also when the measure comes to play, the tool in metric space does not suffice, we have to use a whole different set of tools. 	
\end{remark}

\begin{proof}
	We would like show the uniformly convergence:
\begin{equation}
	\sup |F_n - F| \overset{a.s.	}{\longrightarrow} 0
\end{equation}
 Notice by the empirical probability converges to the real probability, we know:
\begin{equation*}
	F_n(x) \overset{a.s.}{\longrightarrow} F(x), \quad F_n(x-) \overset{a.s.}{\longrightarrow} F(x-), 
\end{equation*}
by the strong law of large number. (Since $(\mathbbm{1}_{(x_i \in (- \infty, x))})_{i \in \mathbb{N}}$ is i.i.d and belongs to $L^2$). We separate the cdf along the y-axis. Define:
\begin{equation}
	x_j := \min \{x : F(x) \leq \frac{j}{N}\}, 
\end{equation}
and
\begin{equation}
	R_n := \max_{j = 1, \dots, N}\{ |F_n(x_j) - F(x_j)| + |F_n(x_j-) - F(x_j-)| \}
\end{equation}
(Since the cdf is only right-continuous we must include the left-limits point)

$\forall x \in (x_{j - 1}, x_j)$ we have:
\begin{equation*}
	F_n(x) \overset{(1)}{\leq} F_n(x_{j}-) \overset{(2)}{\leq} F(x_j-) + R_n \overset{(3)}{\leq} F(x) + R_n + \frac{1}{N} 
\end{equation*}
Where (1) follows by the monotony of $F_n$ and (2) and (3) follow by the definition of $R_n$ and analog we have:
\begin{equation*}
	F_n(x) \geq F_n (x_{j -1}) \geq F(x_{j - 1}) - R_n \geq F(x) - R_n - \frac{1}{N}
\end{equation*}
The characteristic of the uniform convergence is that it does not depends on the selection of $x$. 
Together we have:
\begin{equation}
	\limsup_N \sup_{x \in \Omega}(|F_n - F|) \leq R_n + \frac{1}{N}
\end{equation}
Since $R_n \longrightarrow 0$ by the results above the equation goes to zero as $N \longrightarrow \infty$.
\end{proof}



\begin{itemize}
\color{red}
	\item Shannon's Entropy\\
	Define the probability of a sequence of words:
	\begin{equation*}
		\pi_n = \prod p_i(\omega)
	\end{equation*}
	and the random variable:
	\begin{equation*}
		Y_i(\omega) = -\log(p_i(\omega))
	\end{equation*}
	The Shannon's entropy is defined as:
	\begin{equation*}
		H(p) := -\sum_{e \in E} p_e \log (p_e)) = \mathbb{E}Y
	\end{equation*}
	Apply Law of large number we have $\frac{1}{n} \sum Y_i \longrightarrow H(p)$
\end{itemize}

\subsection*{Aufgabe 3}

\begin{proof}
	$h(S) \longrightarrow h(S_0)$ in probability is proved by the continuous mapping theorem. 
	By Taylor we have:
	\begin{align*}
		\sqrt{n}(h(S_n) - h(s_0) &= \sqrt{n}( h'(s_0)(h(s_n) - h(s_0)) - \frac{h''(\xi)}{2} (h(s_n) - h(s_0))^2)\\
		&= \sqrt{n}h'(s_0)(h(s_n) - h(s_0)) - \frac{C}{\sqrt{n}}[\sqrt{n}h(s_n) - h(s_0)]^2
	\end{align*}
	where the first part converges in distribution to $N(0, h'(s_0)^2 \sigma_g^2 )$ by Slutzky lemma, and the second part converges in probability to 0. 
	
\end{proof}

\begin{remark}
In the multivariate case we have:
\begin{equation*}
	\sqrt{n}( h(s_n) - h(s_0) ) \overset{D}{\longrightarrow} N(0, J(h(s_0)) \Sigma^2 J(h(s_0))^T)
\end{equation*}	
\end{remark}


eigenspace of the ellopsoid 




\subsection*{Aufgabe 4}
\begin{enumerate}
	\item We would like to show:
	\begin{equation}
		f(p) = \mathbb{E}f(\overline{p})
	\end{equation}
	where: $ \overline{p} = \frac{1}{n}\sum_{i = 1}^{n}X_i$
	Denote: $S_n := \sum_{i = 1}^{n}X_i$. We know that $S_n$ is Binomial distributed. Therefore we have:
	\begin{equation}
		\mathbb{E}f(\overline{p} ) = \sum_{k = 0}^{n} f(\frac{k}{n}) \binom{n}{k}p^k(1 - p)^{n - k} = f(p)
	\end{equation}
	\item
	Consider the assumption that f is continuous, therefore we have: $ \forall \epsilon > 0, \exists \delta>0$:
	\begin{equation}
		|f(x) - f(y)| < \epsilon, \forall |x - y| < \delta
	\end{equation}
	We therefore have:
	\begin{equation}
	\begin{aligned}
		|f(\overline{p}) - f(p)| & \leq \epsilon + 2\|f\|_{\infty} \mathbbm{1}_{(|\overline{p} - p| \geq \delta)}\\
		\mathbb{E}|f(\overline{p}) - f(p)| & \leq \epsilon + 2\|f\|_{\infty} \mathbb{P}[|\overline{p} - p| \geq \delta]\\
		& \leq \epsilon + 2\|f\|_{\infty} \frac{Var[\overline{p}]}{\delta^2}\\
		& \leq \epsilon + \frac{\|f\|_{\infty}}{2n\delta^2}
	\end{aligned}
	\end{equation}
	Therefore we have the first convergence result. And by the result of (1) we have:
	\begin{equation}
		\sup_{p \in [0, 1]} |f_n(p) - f(p)|= \sup_p \mathbb{E}|f(\overline{p}) - f(p) |  \long
	\end{equation}
	which tends to zero as n grow large enough.
\end{enumerate}

\includepdf[pages = -]{sheet2.pdf}

\section*{Set 2}

\subsection*{Aufgabe 1}

\subsubsection*{a)}
\begin{proof}
By the properties of gamma distribution, we know for $Y \sim Gamma(\alpha, \beta)$:
\begin{equation*}
	\mathbb{E}_{\alpha, \beta} Y = \frac{\alpha}{\beta}, \quad Var[Y] = \frac{\alpha}{\beta^2} 
\end{equation*}
Define:
\begin{equation*}
	m_1(\alpha, \beta) = \mathbb{E} Y = \frac{\alpha}{\beta}, \quad
	m_2(\alpha, \beta) = \mathbb{E} Y^2
\end{equation*}
\textcolor{red}{Do not mistake the second moment for the variance again!}\\

Now substitute for the empirical measure we have:
\begin{equation*}
\begin{aligned}
	\alpha &= \frac{m_1^2}{m_2 - m_1^2} \\
	\beta &= \frac{m_1}{m_2 - m_1^2}
\end{aligned}
\end{equation*}

The estimator is consistent and therefore root-n normal by exercise 1.2. 

\end{proof}

\textcolor{blue}{
Observe that as $n \longrightarrow \infty$, $\mathbb{E}_{\alpha}\hat{\alpha}_n$ does not converges to $\alpha$. Therefore the estimator is not consistent. \\
Furthermore, as $ \sqrt{n}(\hat{\alpha}_n - \alpha)$ does not converges in distribution to 0, the sequence of estimator does not satisfies the root-n normality.}



\subsubsection*{b)}
\begin{proof}
	

The estimator $\tilde{\theta}$ under the product measure $\mathbb{P} =  \mathbb{P}^{\otimes n}_\theta$ :
\begin{equation*}
	F(a) = \mathbb{P}(\tilde{\theta} \leq a) = \mathbb{P} (x_1 \leq a, \dots, x_n \leq a) = (\frac{a}{\theta})^n 
\end{equation*}
Let $Z := \max_i\{ X_i\}$, we have:
\begin{equation*}
	\begin{aligned}
		\mathbb{P}(Z \leq z) &= (\frac{z}{\theta})^n\\
		f_Z(z) &= \frac{d}{dz}(\frac{z}{\theta})^n = \frac{n}{\theta^n} z^{n - 1} \\
	\end{aligned}
\end{equation*}
Therefore
\begin{equation*}
	\mathbb{E}_\theta[\tilde{\theta}] = \int_0^\theta z \frac{n}{\theta^n} z^{n - 1} dz = \frac{n}{n + 1} \theta
\end{equation*}
The bias:
\begin{equation*}
	\mathbb{E}\tilde{\theta} - \theta = -\frac{1}{n+1} \theta
\end{equation*}
The variance of $\tilde{\theta}:$
\begin{align*}
	\mathbb{E}\tilde{\theta}^2 &= \int_0^\theta z^2 \frac{n}{\theta^n} z^{n-1} dz\\
	 &= \frac{n}{n+2} \theta^2\\
	Var[\tilde{\theta}] &= \mathbb{E} \tilde{\theta}^2 - (\mathbb{E} \tilde{\theta})^2\\
	 &= (\frac{n}{n+2} - (\frac{n}{n+1})^2)\theta^2
\end{align*}
The quadratic risk:
\begin{equation*}
	\mathbb{E} | \tilde{\theta} - \theta |^2 = ( (\frac{1}{n+1})^2 + \frac{n}{n+2} - (\frac{n}{n+1})^2 )\theta^2
\end{equation*}
The estimator is consistent since:
\begin{equation*}
	\mathbb{E}\tilde{\theta}_n - \theta \overset{\mathbb{P}}{\longrightarrow} 0
\end{equation*}
\textcolor{blue}{Observe the root-n normality:}
\begin{equation*}
	\sqrt{n} (\tilde{\theta} - \theta) = \sqrt{n} (\tilde{\theta} - \frac{n}{n + 1} \theta	- \frac{1}{n+1}\theta) = \sqrt{n}(\tilde{\theta} - \mathbb{E}\tilde{\theta}) - \frac{\sqrt{n}}{n+1}\theta
\end{equation*}
\textcolor{blue}{where the first part converges in distribution to a normal distribution and the second part does not assume any randomness and converges to 0. Therefore we have the root-n normality.}\\
\textcolor{red}{The statement above is wrong.This provides the example that, when a moment estimator is not in $C^2$, The consistency doesn't imply root-n normality. \\
	Judging by the rate of convergence. $\max\{Y_i\}$ is a much better estimator than the moment estimator. For the regular family the moment estimator is a good estimator}

\end{proof}


\subsection*{Aufgabe 2}
\subsubsection*{a)}
\begin{equation*}
	\begin{aligned}
		\mathbb{E}_{(\alpha, \sigma^2)} \tilde{\alpha} &= \frac{1}{n}\sum_{i = 1}^n \mathbb{E}Y_i\\
		&= \frac{1}{n} \times n\times \alpha = \alpha\\
		Var_{(\alpha, \sigma^2)}\tilde{\alpha} &= \mathbb{E}(\tilde{\alpha} - \alpha)^2\\
		&= \mathbb{E}(\frac{\sum_{i = 1}^n Y_i}{n} - \alpha)^2\\
		&= \mathbb{E}(\frac{\sum_{i = 1}^n Y_i - n \alpha}{n})^2\\
		&= \frac{1}{n^2} \mathbb{E}(\sum_{i = 1}^n Y_i - n \alpha)^2
	\end{aligned}
\end{equation*}
We know from Stochastic I, that $S_n := \sum_{i =1}^n Y_i \sim \mathcal(n \alpha, n \sigma^2)$, because $Y_i$ i.i.d.. So we have:
\begin{equation*}
	Var_{(\alpha, \sigma)}\tilde{\alpha} = \frac{\sigma^2}{n}
\end{equation*}

\subsubsection*{b)}
\textcolor{red}{The following proof is inacurate. It does not consider the multidimensional case.}
\begin{proof}
	The goal is to establish a confident ellipsoid (In this case a confidence interval) for $\alpha$ at the level of $c$. By definition:
\begin{equation}
	E(z_c) = \{ \alpha \in \mathbb{R}\mid \frac{\sqrt{n}}{\sigma} \mid \tilde{\alpha} - \alpha \mid \leq z_c\}
\end{equation}
Notice that the variance here is unknown. We replace the true $\sigma$ by the consistent estimate $\tilde{\sigma}$ and get the data-driven estimate:
\begin{equation*}
	E^\circ (z_c) = \{\alpha \in \mathbb{R} \mid -\frac{z_c \sigma}{\sqrt{n}} \leq \alpha \leq \frac{z_c\sigma}{\sqrt{n}}\}
\end{equation*}
\end{proof}

\begin{proof}
	by the exercise 1.2. In order to calculate the confidence ellipsoid, we need to know: $\sigma_{g(x)}$, $g(x)$ in this case is $(y, y^2)$, and $J_{m^{-1}}$. We get the asymptotic distribution:
	\begin{equation*}
		\sqrt{n}(\theta^* - m^{-1} (g(y)) ) \longrightarrow N(0, J_{m^{-1}} \sigma^2_{g(x)} J_{m^{-1}}^T)
	\end{equation*}
	define $J_{m^{-1}} \sigma^2_{g(x)} J_{m^{-1}}^T := \Sigma$. We have:
	\begin{equation*}
		\Sigma^{-\frac{1}{2}}(\theta^* - \tilde{\theta}) \longrightarrow N(0, I_p)
	\end{equation*}
	where $p$ is the dimension of the parameter. 
	The confidence ellipsoid has the form:
	\begin{equation*}
		\{\theta^* | (\theta^* - \tilde{\theta})\Sigma(\theta^* -\tilde{\theta)} \leq \frac{q_{\chi_p}(\alpha)}{n}\}
	\end{equation*}
\end{proof}


\subsection*{Aufgabe 3}
Proof the multivariate bias-variance decomposition:
\begin{equation*}
	R(\theta^*, \tilde{\theta}) = tr(Var(\tilde{\theta})) + \| bias(\theta^*, \tilde{\theta})\|^2
\end{equation*}


\subsection*{Aufgabe 4(Aspt CFs)}
Suppose ($\tilde{\theta}_n$) is a root-n normal sequence of estimator of the true parameter $\theta^*$. $\sigma$ is continuous on $\Theta$. To show the confidence sets motivated by central limits theory satisfies the purpose.
\begin{proof}
	
\end{proof}

\includepdf[pages = -]{sheet3.pdf}
\section*{Set 3}
\begin{itemize}
\color{red}
\item Equality in the Jensen'Inequality: The equality holds if and only if the convex function is affine or the argument is constant.
\end{itemize}


\subsection*{Aufgabe 1}
The KL-divergence usually is considered as loss function. 
\begin{proof}.
\begin{enumerate}
	\item Positive definite:\\
	First we show that $\mathbb{D}_f(\mathbb{P}, \mathbb{Q}) \geq 0$. Apply Jensen Inequality, since $f$ is convex, we have:
	\begin{equation*}
		\mathcal{D}_f(\mathbb{P}, \mathbb{Q}) = \mathbb{E}_Qf(\frac{dP}{dQ}) \geq f(\mathbb{E}_Q\frac{dP}{dQ}) = f(1) = 0
	\end{equation*}
	Then we show: $ \mathbb{P}, \mathbb{Q}$ are equivalent $\rightleftarrows$ $\mathcal{D}_f(\frac{dP}{dQ}) = 0$\\
	"$\Rightarrow$" $\mathbb{P}, \mathbb{Q}$ equivalent $\Rightarrow$ $\frac{dP}{dQ} = 1, a.s.$ $\Rightarrow$ $f(\frac{dP}{dQ}) = 0, a.s.$ $\Rightarrow$ $\mathcal{D}_f(\mathbb{P}, \mathbb{Q}) = 0$\\
	"$\Leftarrow"$, we know in Jensen's inequality, the equality holds if and only if $f$ is affine or the argument is constant. Since $f$ is \textcolor{red}{strictly} convex, $\frac{d\mathbb{P}}{d\mathbb{Q}}$ must be constant. They are both probability measure so it could only be 1. 
	\item Convexity:
	\begin{itemize}
		\item to show $\alpha \mathbb{P} + (1-\alpha) \mathbb{P} \ll \alpha \mathbb{Q} + (1 - \alpha) \mathbb{Q}$.\\
		let $A \in \mathcal{F}$ a null set in the new measure. By definition we have:
		\begin{align*}
			\forall \alpha, \alpha \mathbb{Q}_1 [A] + (1 - \alpha)\mathbb{Q}_2 [A] = 0
		\end{align*} 
		Let $\alpha = 0, \alpha = 1$ respectively. We have $A$ is a null set in both $\mathbb{Q}_1, \mathbb{Q}_2$. \\
		$\Rightarrow$ $A$ is a null set in $\mathbb{P}_1, \mathbb{P}_2$. \\
		$\Rightarrow$ $A$ is a null set in the convex combination $\alpha\mathbb{P}_1 + (1 -\alpha)\mathbb{P}_2$.\\
		We have shown the absolute continuity.
		\item convexity: By showing the determinant of the Hessian matrix of the function $G: (x, y) \mapsto yf(\frac{x}{y})$ is non-negative. Compute the Jacobi-Matrix:
		\begin{equation*}
			J_G = (f'(\frac{x}{y}), f(\frac{e}{y}) - \frac{x}{y} - \frac{x}{y}f'(\frac{x}{y}))^\top
		\end{equation*}
		The corresponding Hessian Matrix is:
		\begin{equation*}
		H_G =
		\begin{bmatrix}
			\frac{1}{y}f''(\frac{x}{y}) & xf''(\frac{x}{y})\\
			\frac{x}{y^2}f''(\frac{x}{y}) & \frac{x^2}{y^3}f''(\frac{x}{y})
		\end{bmatrix}
		\end{equation*}
		Therefore $\det_{H_G} \geq 0$
	\end{itemize}
	\item Partition Property\\
		followed by Jensen and convexity.	
\end{enumerate}
\end{proof}

\subsection*{Aufgabe 2}

\subsubsection*{1)}
To show the equivalence of the variation distance:\\
Since the $f-divergence$ is essentially $Q-Integral$ and integral is defined by supremum of step function. We just need to $\mathbb{E}_Q f$ could be written as a step function that conforms with $\sup_{A \in \mathcal{F}} \mid \mathbb{P}(A) - \mathbb{Q}(A) \mid$.
\begin{equation*}
\begin{aligned}
	\mathbb{E}_Q f &= \int_\Omega \frac{1}{2} \mid \frac{d\mathbb{P} - d\mathbb{Q}}{d \mathbb{Q}} \mid d\mathbb{Q}\\
	&= \frac{1}{2} \int_\Omega \mid d\mathbb{P} - d\mathbb{Q} \mid d \lambda\\
	&= \int_{dP \geq dQ} d\mathbb{P} - d\mathbb{Q}  d \lambda\\
	&= \sup_{A \subset \Omega m.b.} | \mathbb{P}(A) - \mathbb{Q}(A) |
\end{aligned}
\end{equation*}

\subsection*{Aufgabe 4}
Calculation of the Fischer Information. Remember we have 2 ways of calculating the Fischer information.



\includepdf[pages = -]{sheet4.pdf}
\section*{Set 4}
\subsection*{Aufgabe 1}
\subsubsection*{a)}

\begin{equation*}
	\mathbb{F} = - \int l''
\end{equation*}

Since both $(Poisson(\theta), \theta > 0)$,  $(\mathcal{N}(\theta, \sigma^2)$ for some fixed $\sigma^2$ are natural exponential family, we know by the theorem in the lecture that the empirical mean is a R-efficient estimator.\\


\subsection*{Aufgabe 2}
\subsubsection*{a)}
\begin{proof}
	Suppose $\mathbb{P} \ll \mu$, the Fisher information matrix has the form:
	\begin{align*}
		F(\theta)_{ij} &= \mathbb{E}_\theta \nabla_i l(\theta, Y) \cdot \nabla_j l(\theta, Y)\\
		&= \int \nabla_i l(\theta, Y) \cdot \nabla_j l(\theta, Y) \frac{d\mathbb{P}}{d\mu} d\mu
	\end{align*}
	This integration is independent of the dominating measure $\mu$.\\
	We should further show that the gradient of the log-likelihood function is independent of the
dominating measure. Suppose $\mathbb{P} \ll \mu_1, \mathbb{P} \ll \mu_2$, and define:$\mu_3 = \mu_1 + \mu_2$. It is trivial to show that $\mu_1, \mu_2$ is absolutely continuous w.r.t. $\mu_3$. Observe now:
\begin{align*}
	\nabla_i \log(\frac{d\mathbb{P}}{d\mu_3}) &= \nabla_i \log(\frac{d \mathbb{P}}{d \mu_1} \frac{d\mu_1}{d\mu_3}) \\
	&= (\frac{d\mathbb{P}}{d\mu_1} \frac{d\mu_1}{d\mu_3})^{-1} \nabla_i(  \frac{d\mathbb{P}}{d\mu_1}\frac{d\mu_1}{d\mu_3})\\
	&= \nabla_i \log(\frac{d\mathbb{P}}{d\mu_1})
\end{align*}
Analog we can show the $\nabla_i \log (\frac{d\mathbb{P}}{d\mu_2}) = \nabla_i \log (\frac{d\mathbb{P}}{d\mu_3})$.\\
Therefore we see the gradient of the log-likelihood function is not dependant on the dominating measure. 
\end{proof}

\subsubsection*{b)}
\begin{proof}
\begin{align*}
	(\mathbb{F}_{\theta^{\otimes n}}(\theta))_{ij} &= \mathbb{E}_{\theta^{\otimes n}} \nabla_i \sum_{i = i}^n l(\theta, y_i) \cdot \nabla_j \sum_{i = i}^n l(\theta, y_i)\\
	(independence)&= \mathbb{E}_{\theta^{\otimes n}} \sum_{k = 1}^n \nabla_i l(\theta, y_i) \nabla_j l(\theta, y_i) \\
	(Fubini)&= \mathbb{E}_\theta \nabla_i l(\theta, y_1) \nabla_j l(\theta, y_1)\\
	&= (\mathbb{F}(\theta))_{ij}
\end{align*}
\end{proof}
need to be refined.


\subsubsection*{c)}
$\mathbb{F}(\theta) = Var \nabla l (Y,\theta)$
\begin{proof}
	We show that $\mathbb{E} \nabla l(Y, \theta) = 0$, then by the definition of Fischer Information Matrix we could have the result.
	By the definition of regular statistical model:
	\begin{align*}
		\mathbb{E}_\theta \nabla l(Y, \theta) &= \mathbb{E}_\theta\nabla\frac{d\mathbb{P}_\theta}{d\mu} \quad \text{(Definition of log-likelihood)}\\
		 &= \nabla \mathbb{E}_\theta \frac{d\mathbb{P}_\theta}{d\mu} (Y, \theta) \quad \text{(Definition of regular model)}\\
		 & = \nabla 1 = 0
	\end{align*}
\end{proof}
.\\
$\mathbb{F}(\theta) = -\mathbb{E} Hess(l(Y, \theta))$
\begin{proof}
\end{proof}


\subsubsection*{d)}

\begin{proof}
	Define $m(\theta): = \mathbb{E}_\theta l(\theta, Y)$, then we have:
	\begin{align*}
		KL(\theta, \theta') &= m(\theta) - m(\theta')\\
		&= \langle \nabla m(\theta), \theta - \theta' \rangle + \frac{1}{2} (\theta - \theta ')^\top Hess(\xi) (\theta - \theta') \quad \xi \in (\theta, \theta') \\
		&= \langle \nabla \mathbb{E}_\theta l(\theta, Y), \theta - \theta' \rangle - \frac{1}{2} (\theta - \theta ')^\top \mathbb{F}(\xi) (\theta - \theta')\\
		&= (\theta - \theta ')^\top \mathbb{F}(\xi) (\theta - \theta')
	\end{align*}
\end{proof}


\includepdf[pages = -]{sheet5.pdf}
\section*{Set 5}
\subsection*{Aufgabe 1}
\subsubsection*{1)}


\subsubsection*{2)}

\subsubsection*{3)}
\begin{proof}
	\begin{equation*}
		\mathbb{E} \psi(\theta, y) = \int_{-\infty}^\theta (y - \theta)(\alpha - 1) d\mathbb{P}_\theta(y)
		 + \int_\theta^\infty \alpha(y - \theta) d\mathbb{P}(y)
	\end{equation*}
\end{proof}

\subsubsection*{4)}
\begin{proof}
	Assumed the errors are Laplacian distributed:
	\begin{equation*}
		\epsilon \sim p(x) = \frac{1}{2}\exp(\frac{-|x|}{2})
	\end{equation*}
\end{proof}



\subsection*{Aufgabe 3}

\begin{proof}
	the consistent estimator



	We know the MLE of an EFn is $\frac{S}{n}$, which is R-efficient. (By the definition of natural exponential family it is obviously unbiased). Therefore we can construct the confidence interval by normal approximation:
	\begin{equation*}
		E = [\tilde{\theta}_n - \frac{\sigma}{\sqrt{n}}z_\alpha, \tilde{\theta}_n + \frac{\sigma}{\sqrt{n}}z_\alpha]
	\end{equation*}
	This interval is alway symmetric and the probability, that the true parameter does not lie in the confidence set, is asymptotically approximating $\alpha$.
\end{proof}



\includepdf[pages = -]{sheet6.pdf}
\section*{Set 6}

\subsection*{Aufgabe 1}
\subsubsection*{a)}
Given a quasi linear model, calculate the $L(\theta)$.
\begin{proof}
	\begin{equation*}
		Y - \psi^T \theta \sim N(0, \Sigma_0)
	\end{equation*}
\end{proof}


\subsubsection*{b)}
\begin{proof}
	The formula is given by the normal equation:
	\begin{equation*}
		\tilde{\theta} = (\psi \Sigma^{-1} \psi^T)^{-1} \psi \Sigma^{-1} Y
	\end{equation*}
	The invertability is a necessary condition. 
\end{proof}

\subsubsection*{c)}
$Var(\tilde{\theta})$ could be shown by the standard formula.

\subsubsection*{d)}
\begin{proof}
	
\end{proof}







\section*{Set 7}

\subsection*{Aufgabe 1}[Fisher Information Inequality]
\begin{proof}
Observe that $KL(\theta_1, \theta_2)$ admits the minimum at $\theta_1$.
	\begin{align*}
		KL^{\frac{1}{2}} (\theta_1, \theta_3) &\leq (\theta_1 - \theta_3)\sqrt{c_2}\\
		&= \sqrt{c_2} (\theta_1 - \theta_2 + \theta_2 - \theta_3)\\
		&\leq \sqrt{\frac{c_2}{c_1}} (KL^{\frac{1}{2}}(\theta_1, \theta_2) + KL^{\frac{1}{2}}(\theta_2, \theta_3))
	\end{align*}
\end{proof}



\subsection*{Aufgabe 2}


\section{Technique}
\subsection{Taylor expansion}

\section{Prong to error}
\subsection{Scalar transformation for matrix}
\begin{equation*}
	\langle A, CB \rangle = \langle AC^T, B\rangle
\end{equation*}
the side of tranformation

\subsection{variance of rv multiplied by matrix}
\begin{equation*}
	Var(A\mathbb{Y}) = A Var(\mathbb{Y}) A^T
\end{equation*}

\section{Techniques}
\subsection{Rewrite the characteristic funtions}

\end{document}
